\begin{textsample}{POS Dim 3 – vem\_ed\_21 – Score 2.00 – vem\_ed\_21/t108.txt}  \label{ex:f3_pos_110}
continuação Apoio ao matchmaking entre parceiros, treinamento docente, orientação e suporte, reconhecimento de professores e alunos estão entre as recomendações para institucionalizar as iniciativas COIL.
A segunda comunicação de Ana Carolina foi"\textbf{Regional} campuses in the lead: Institutionalizing COIL across multi-campus public systems"(Campi \textbf{regionais} na liderança: institucionalizando COIL em sistemas públicos multicampi), com Dan Nolan (University of Minnesota System / EUA) e Natalia Dyba (University of Washington Bothell / EUA).
Os autores relataram os esforços de \textbf{institucionalização} de campi \textbf{regionais} em sistemas públicos no Brasil e nos Estados Unidos.
Destacaram que a organização dos Intercâmbios Virtuais nesse contexto requer alinhamento com metas interinstitucionais mais amplas e uma noção dos pontos fortes, da preparação e da capacidade do corpo docente, da equipe e da liderança.

% matched lemmas: institucionalização, regional
\end{textsample}
